% =====================================================================
%  Offline On-Policy Distillation Breaks Under Multi-Turn Distribution Shift
%  A controlled, GPU-free study.
%
%  OVERLEAF UPLOAD INSTRUCTIONS  (IMPORTANT: read the compiler line)
%  ----------------------------
%  Upload this file (report.tex) together with the figures:
%      architecture.png
%      fig_main.png  fig_coverage.png  fig_deploy.png  fig_representation.png
%      fig_pareto.png  fig_sweeps.png
%  (they live in ../results/ on disk; upload them flat, next to this .tex).
%
%  >>> SET THE COMPILER TO XeLaTeX <<<
%  In Overleaf: Menu (top-left)  ->  Settings  ->  Compiler  ->  XeLaTeX.
%  The ctex package below needs XeLaTeX for full Chinese support. With it you
%  may type Chinese ANYWHERE (section titles, body, captions) and edit freely.
%  Do NOT compile with pdfLaTeX -- Chinese will fail.
% =====================================================================
\documentclass[11pt]{article}

% Compiler guard: MUST use XeLaTeX
\RequirePackage{ifxetex}
\ifxetex\else
  \errmessage{%
    ^^J========================================^^J
    COMPILER ERROR: You must use XeLaTeX!^^J
    In Overleaf: Menu -> Settings -> Compiler -> XeLaTeX^^J
    ========================================^^J}
\fi

\usepackage{ctex}              % full Chinese support; requires XeLaTeX
\usepackage[margin=1in]{geometry}
\usepackage{booktabs}
\usepackage{graphicx}
\usepackage{amsmath}
\usepackage{amssymb}
\usepackage{xcolor}
\usepackage{array}
\usepackage{caption}
\usepackage{enumitem}
\usepackage{fontawesome5}      % \faGithub icon for the repo link
\usepackage[colorlinks=true,linkcolor=blue,urlcolor=blue,citecolor=blue]{hyperref}
\usepackage{url}

% Compact, paper-style spacing
\captionsetup{font=small,labelfont=bf}
\setlength{\parskip}{0.4em}
\setlength{\parindent}{0pt}

% Shorthand for inline code
\newcommand{\code}[1]{\texttt{#1}}
% Inline CJK -- with ctex (XeLaTeX) Chinese works everywhere, so this is a
% transparent pass-through. Kept so existing \cn{...} calls still compile.
\newcommand{\cn}[1]{#1}
% Section question tag, right-aligned and light
\newcommand{\qtag}[1]{\hfill{\normalfont\small\color{gray}[Question (#1)]}}

\title{\textbf{Offline On-Policy Distillation Breaks Under\\ Multi-Turn Distribution Shift}\\[0.3em]
\large A Controlled, GPU-Free Study}

\author{Jingyi Du \,(\cn{杜兢漪})\\[0.2em]
\normalsize First-year Undergraduate, Joint Institute (\cn{密院大一本科生})\\
\normalsize Shanghai Jiao Tong University\\
\normalsize \href{mailto:productivity25@sjtu.edu.cn}{productivity25@sjtu.edu.cn}}
\date{}

\begin{document}
\maketitle

\begin{center}
\small Code: \faGithub\;\url{https://github.com/Dujingyi2006/Offline-OPD-under-Agentic-Shift}
\end{center}

% ---------------------------------------------------------------------
\paragraph{Why I picked this question.}
I want to join the Industrial Foundation Group as a research intern. Foundation
model training is the center of gravity in AI right now, and I want to be part
of it. I have no model training experience yet, and I have never touched a GPU
cluster, but I am deeply interested in how these systems learn, where they
break, and how to fix them. I know Professor [Name] runs an 800-A100 cluster,
one of the largest academic training setups in the field. That is exactly the
environment I want to learn in. This report is my way of showing I can think
carefully about training dynamics, isolate failure modes, and propose testable
fixes, even when I have to build the entire experiment on a CPU because I do
not yet have access to the real thing. I chose the offline distillation question
because it sits at the intersection of efficiency (training without environment
access, the thing that scales) and the multi-turn agentic setting (the thing
foundation models are increasingly asked to do). If I can show the method breaks
under controlled conditions and explain why, that is a step toward understanding
what would break at scale.

\paragraph{Abstract.}
Lightning OPD offers the on-policy signal of online distillation at the price of
supervised fine-tuning, by moving the teacher advantage off-policy and removing the
environment from the training loop. Its theory rests on two assumptions, support
coverage and teacher consistency, and argues that under consistency any residual
gap between the online and offline updates is governed by policy drift and is otherwise
capacity-limited. We test this central claim in the setting the paper itself flags as
open, a multi-turn agentic task. We build a small fault-diagnosis environment that
runs on a CPU, with no GPU or large language model, in which an early tool error pushes
the agent onto a recovery branch the clean demonstrations never visit. A single tabular
oracle writes the demonstrations and supplies every teacher target, so teacher
consistency holds by construction, and the recovery action is representable in the
student, so capacity is not the bottleneck. Across five seeds, Lightning OPD lands
at or below the supervised floor ($0.632$ vs $0.699$ success) with roughly ten
times the variance, while online RL and online OPD reach $1.000$. The surviving gap is
therefore neither a consistency bias nor a capacity limit. We trace it to a third
cause, coverage of the distilled signal. Off-support, the sampled-action advantage
is a one-sample estimator drawn from the reference policy, which places near-zero mass
on the recovery action, so the recovery term is not high-variance but systematically
absent from the estimator's support. Three ablations confirm the mechanism: plain
Lightning OPD is flat across collection coverage, tracks supervised fine-tuning under
deployment shift, and collapses to the supervised floor once parameter generalization
is removed. We then change one line, replacing the sampled action by the closed-form
expectation under the full teacher distribution, and the same student reaches
$1.000$ at the same one-time collection budget, with no environment access and no live
teacher, matching the full-distribution offline method of Rang et al.\
(arXiv:2509.26497) cell for cell. Two caveats remain. This $1.000$ is
argmax-aligned and falls to $\approx 0.87$ under sampled evaluation, which a
branch-aware replay reweighting then lifts to $0.95$, still at zero environment cost,
while an uncertainty-triggered teacher query that also closes the gap does so only
by spending the environment access the method exists to avoid. The operative variable
is not whether one touches the environment, but what one distills.

\hrulefill

\newpage
% =====================================================================
\thispagestyle{empty}
\begin{center}
{\Large\textbf{Contents}}\\[0.3em]
{\small 教授布置的七道题分别在以下章节作答。}
\end{center}
\vspace{1.5em}

\begin{center}
\renewcommand{\arraystretch}{1.6}
\begin{tabular}{c p{10.5cm} c}
\toprule
\textbf{题号} & \textbf{教授要求（原题）} & \textbf{对应章节} \\
\midrule
(i)   & 阅读 Lightning OPD 论文以及其 codebase，理解 SFT、online RL、online OPD、offline OPD 的区别。 & \S\ref{sec:methods} \\
(ii)  & Lightning OPD 在 offline 情况下做了哪些假设？这些假设在所有类型的数据上都成立吗？ & \S\ref{sec:assumptions} \\
(iii) & 设计一个轻量级的多轮 agentic toy task（无需 GPU 与 LLM）：包含多轮决策、工具式动作、错误分支、最终答案奖励，且早期错误会污染后续状态。 & \S\ref{sec:task} \\
(iv)  & 实现并对比至少三种方法（SFT / online RL / offline OPD，可选 online OPD 作为上界）：成功率、平均奖励、轨迹长度、off-support 状态比例。 & \S\ref{sec:compare} \\
(v)   & Lightning OPD 在多轮 agentic 任务中有哪些问题？offline OPD 在什么情况下有效、什么情况下失效（分布偏移、状态覆盖不足、错误轨迹累积）？ & \S\ref{sec:why} \\
(vi)  & 与 \href{https://arxiv.org/pdf/2509.26497}{arXiv:2509.26497} 中的离线策略相比有哪些区别？ & \S\ref{sec:rang} \\
(vii) & 针对观察到的某一个局限，提出一个代码/方法层面的改进，并做一个简单的消融实验。 & \S\ref{sec:patch} \\
\bottomrule
\end{tabular}
\end{center}
\vspace{1.5em}

\begin{center}
\small This report keeps a paper-style layout. Each section heading also quotes the
assigned question it answers, so the two readings line up one to one.
\end{center}

\newpage

% =====================================================================
\section{阅读 Lightning OPD 论文以及其 codebase，理解 SFT、online RL、online OPD、offline OPD 的区别 \qtag{i}}
\label{sec:methods}
% =====================================================================

I read the Lightning OPD paper (\code{t2.pdf}) and its codebase. Four training
recipes appear. They differ on two axes: does the environment run inside the
training loop, and does a live teacher run inside it. Lightning OPD's whole pitch
is to answer ``no'' to both, yet keep an on-policy-style signal.

The advantage signal is the same object in online and offline OPD. For a state
$s_t$ and action $a_t$,
\[
A_t = \log \pi_T(a_t \mid s_t) - \log \pi_\theta(a_t \mid s_t)
\qquad \text{(stop-grad, clipped).}
\]
Online OPD computes it on states the current student visits. Lightning
OPD computes it once, on rollouts sampled from a frozen reference policy
$\pi_{\mathrm{ref}}$, and reuses that frozen signal for every gradient step.

\begin{table}[h]
\centering
\small
\begin{tabular}{l c c l r}
\toprule
\textbf{method} & \textbf{env in loop?} & \textbf{live teacher?} & \textbf{what it optimizes} & \textbf{env steps} \\
\midrule
SFT          & no  & no  & cross-entropy on fixed demos                         & 0 \\
online RL    & yes & no  & policy gradient on env reward                        & $\sim$155k \\
online OPD   & yes & yes & advantage $A_t$ on freshly visited states            & $\sim$30k \\
offline OPD  & no  & no  & frozen advantage $A_t$ on $\pi_{\mathrm{ref}}$ rollouts & $\sim$635 \\
\bottomrule
\end{tabular}
\caption{The four recipes. Env steps are the measured training-time cost in this
study (\S\ref{sec:patch}). Offline OPD's 635 is a one-time collection that
350 training iterations then reuse with zero further \code{env.step}.}
\label{tab:methods}
\end{table}

Read the table top to bottom and the trade is clear. SFT is the cheapest and the
most fragile: it copies demonstrations and never sees its own mistakes. Online RL
and online OPD are correct but expensive, because they pay environment steps on
every gradient step. Lightning OPD is the interesting case. It claims the
on-policy signal of OPD at the price tag of SFT. The rest of this report tests
whether that claim survives contact with a multi-turn task.
% =====================================================================
\section{Lightning OPD 在 offline 情况下做了哪些假设？这些假设在所有类型的数据上都成立吗？ \qtag{ii}}
\label{sec:assumptions}
% =====================================================================

To move the advantage off-policy, Lightning OPD leans on two assumptions. Support
coverage (Asm 3.2) requires the frozen $\pi_{\mathrm{ref}}$ to visit every state the
student will later visit; where $\pi_{\mathrm{ref}}$ never went, there is no advantage
to read. Teacher consistency (Asm 3.4) requires the model that wrote the SFT demos and
the model that supplies $\log\pi_T$ to be the same, so the mismatch variance is
$\sigma_\Delta^2 = 0$.

The paper's theory then rests on these. Theorem 3.5 bounds the
online--offline gap by drift alone,
\[
\bigl\|\nabla J_{\mathrm{on}} - \nabla J_{\mathrm{off}}\bigr\|_2
\;\le\; G\,\sigma_A \sqrt{\chi^2(\pi_\theta \,\|\, \pi_{\mathrm{ref}})},
\]
and the gap is exactly zero at initialization $\pi_\theta=\pi_{\mathrm{ref}}$.
Theorem 3.6 gives a shared fixed point when the teacher is representable.
Just after it, the paper goes one step further: when the teacher is not
representable the offline update still ``stays close to the online one as long as
drift remains small, suggesting comparable capacity-limited behavior in practice.''

Do these hold on all data? No. Single-turn data hides the gap; multi-turn
data exposes it. On a single-turn task the student acts once, so it cannot drift
into states $\pi_{\mathrm{ref}}$ never saw. Coverage holds for free and the Thm 3.5
bound is tight. A multi-turn agent is different. One wrong tool call lands it in a
state off the reference distribution. There two things fail at once: the
precomputed advantage is undefined on the action actually sampled, and the
recovery skill was never on the golden path to be sampled in the first place.
Off-support the drift $\chi^2(\pi_\theta\|\pi_{\mathrm{ref}})$ blows up, so the
Thm 3.5 bound becomes empty. The guarantee does not break loudly. It simply stops
saying anything.

The claim I attack. Beating on Asm 3.2 alone is cheap, since the authors
themselves flag multi-turn agents as open in Appendix E. So I aim higher, at the
capacity reading: that under teacher consistency, any surviving gap is a
representability problem, and adding capacity closes it. I design the experiment
to let me test that claim on its own terms. Consistency holds by construction,
because a single tabular oracle teacher writes the SFT demos and supplies every
$\log\pi_T$ target, so $\sigma_\Delta^2 = 0$ exactly. I sit inside the regime where
the paper's guarantee is supposed to bite, which kills the easy referee escape,
``you broke consistency.'' Capacity is not the bottleneck either: the recovery action
is representable in the student's class, and the same linear student, fed the full
teacher distribution (patch 2, \S\ref{sec:patch}), reaches 1.000 at zero extra
env cost. The capacity to recover is present; the sampled-action signal just
never points at it.

With consistency satisfied and capacity sufficient, the capacity reading predicts
near-online behavior. \S\ref{sec:compare} proves it wrong: the method lands at
or below SFT. The surviving gap is therefore neither a consistency bias nor a
capacity limit. It is a coverage-of-the-distilled-signal failure.

% =====================================================================
\section{设计一个轻量级的多轮 agentic toy task（无需 GPU 与 LLM）：包含多轮决策、工具式动作、错误分支、最终答案奖励，且早期错误会污染后续状态 \qtag{iii}}
\label{sec:task}
% =====================================================================

I designed a fault-diagnosis episode. It is small, runs on a CPU, and meets every
requirement in the brief: multi-turn decisions, tool-style actions, an error
branch, a final-answer reward, and early errors that poison later states.

The expert walks a thin golden path. It calls the discriminative tool,
reads the result, then commits the matching fix. Deployment adds tool noise. A
discriminative tool can return a false negative and rule out the true
fault. That pushes the agent onto an error branch. The only escape is
\code{ask\_followup}, a recovery action that never appears on the clean golden
path, so it is absent from any frozen-$\pi_{\mathrm{ref}}$ dataset.

The student is a feature-based linear-softmax policy with shared weights.
This matters. Shared weights mean the student generalizes its on-path
behaviour into unseen states, the way a real model does. There it fails
deterministically, rather than escaping by luck. The error compounds: one bad
tool call early changes every state that follows.

Figure~\ref{fig:arch} lays out the full setup end to end: the offline collection
that freezes the dataset, the advantage-weighted gradient the student trains on, the
multi-turn task it is deployed into, and the precise point where the contract breaks.
The advantage is read at a single action $a\sim\pi_{\mathrm{ref}}$, so the recovery
action on the error branch is never distilled. That is the coverage-of-the-distilled-signal
gap this report measures, not an env-access gap.

\begin{figure}[h]
\centering
\includegraphics[width=\textwidth]{architecture.png}
\caption{Offline On-Policy Distillation under agentic shift: frozen data, no
environment in the loop, no live teacher. The frozen reference $\pi_{\mathrm{ref}}$
records one action per state; the advantage $A=\log\pi_T-\log\pi_\theta$ is read at
that single action. When deploy noise pushes the student onto the error branch, the
recovery action \code{ask\_followup} ($\pi_{\mathrm{ref}}\approx 10^{-8}$) carries no
gradient, so the support-coverage assumption fails as an importance-sampling condition
and success lands at 0.63, below the SFT floor.}
\label{fig:arch}
\end{figure}

% =====================================================================
\section{实现并对比至少三种方法（SFT / online RL / offline OPD，可选 online OPD 作为上界）：成功率、平均奖励、轨迹长度、off-support 状态比例 \qtag{iv}}
\label{sec:compare}
% =====================================================================

I implemented all four required methods plus online OPD as an upper bound, and the
Rang-style offline KD as the full-distribution comparison. Fixed operating point:
\code{collect\_noise = 0.05}, \code{deploy\_noise = 0.30}. The brief asks for four
metrics: success rate, average reward, trajectory length, and off-support state
ratio. Table~\ref{tab:metrics} reports all four, mean $\pm$ std over 5 seeds
(\code{results\_main.csv}, Figure~\ref{fig:main}).

\begin{table}[h]
\centering
\small
\begin{tabular}{l c c c c}
\toprule
\textbf{method} & \textbf{success} & \textbf{avg reward} & \textbf{traj.\ length} & \textbf{off-support ratio} \\
\midrule
teacher (oracle)               & $1.000 \pm .000$ & $0.928 \pm .001$ & $3.61 \pm .03$ & $0.078$ \\
online RL                      & $1.000 \pm .000$ & $0.928 \pm .001$ & $3.61 \pm .04$ & $0.198$ \\
online OPD (upper bound)       & $1.000 \pm .000$ & $0.928 \pm .001$ & $3.61 \pm .03$ & $0.078$ \\
\textbf{Rang et al.\ offline KD} & $\mathbf{1.000 \pm .000}$ & $0.928 \pm .001$ & $3.61 \pm .03$ & $0.078$ \\
\textbf{offline OPD + branch replay (patch 3)} & $\mathbf{1.000 \pm .000}$ & $0.928 \pm .001$ & $3.61 \pm .03$ & $0.078$ \\
\textbf{offline OPD + full-dist adv (patch 2)} & $\mathbf{1.000 \pm .000}$ & $0.928 \pm .001$ & $3.61 \pm .03$ & $0.078$ \\
offline OPD + active query (patch 1) & $0.880 \pm .13$ & $0.55 \pm .24$ & $3.80 \pm .21$ & $0.134$ \\
SFT                            & $0.699 \pm .014$ & $0.319 \pm .028$ & $3.90 \pm .04$ & $0.000$ \\
\textbf{Lightning OPD}         & $\mathbf{0.632 \pm .269}$ & $0.183 \pm .552$ & $4.01 \pm .69$ & $0.147$ \\
\bottomrule
\end{tabular}
\caption{Four-metric comparison at the fixed operating point. Off-support ratio is
the fraction of visited states absent from the frozen $\pi_{\mathrm{ref}}$ dataset.}
\label{tab:metrics}
\end{table}

Read the four metrics together and the failure has a fingerprint.

On success, Lightning OPD lands at 0.632, below the 0.699 SFT
floor. It buys nothing over plain imitation here. Its std is 0.269, an order
of magnitude above SFT's 0.014. That variance is the tell: whether a seed's
frozen dataset happened to brush the error branch decides the run. Average reward
tells the same story, sharper. Lightning OPD averages 0.183 against SFT's 0.319;
the frozen advantage actively drags reward below imitation. Trajectory length adds
the third symptom: Lightning OPD's episodes are the longest (4.01), so it wanders,
while the successful methods finish in $\sim$3.6 steps like the teacher. The
off-support ratio is the diagnostic axis. SFT stays rigidly on-support (0.000),
failing by stalling rather than by wandering. Lightning OPD is the highest among the
no-env methods (0.147): it leaves the covered region and fails there. The winners sit
at the teacher's 0.078, because they recover and rejoin near-golden trajectories.
Online RL's 0.198 is benign, since it trained on those states, so visiting them is fine.

Three methods reach 1.000 from the same frozen, no-env-access budget:
Rang's full-distribution KD, and, decisively, Lightning OPD's own gradient
with one line changed (patch 2, \S\ref{sec:patch}). The variable is what you
distill, not whether you touch the environment.

% =====================================================================
\section{Lightning OPD 在多轮 agentic 任务中有哪些问题？offline OPD 在什么情况下有效、什么情况下失效（分布偏移、状态覆盖不足、错误轨迹累积）？ \qtag{v}}
\label{sec:why}
% =====================================================================

\subsection{The mechanism, from the case studies}

\code{case\_studies.txt} shows a forced error-branch episode (true fault
\code{db}, a discriminative tool gives a wrong reading at step 1). After the false
negative, Lightning OPD cycles uselessly through \code{inspect\_log} $\to$
\code{check\_config} $\to$ \code{query\_db} (twice) and times out at $-1.12$. The
teacher hits the same false negative, plays \code{ask\_followup} at the all-against
state, recovers the belief, and commits \code{final\_db} for $+0.98$, so the recovery
skill exists and OPD just never saw it. Active query (patch 1) queries the teacher at
the uncertain state, gets the recovery distribution, and climbs to 0.88, but at the
cost of 159 live annotations and mid-training env rollouts.

This is a coverage failure, not an optimization failure. The advantage
$\log\pi_T - \log\pi_\theta$ is defined only where the frozen rollouts went. The
error branch is off-support by construction. Lightning OPD has no term that pulls
the student toward the teacher there.

\subsection{Ablation A --- collection coverage}
\code{results\_coverage.csv}, Figure~\ref{fig:coverage} (and panel (a) of
Figure~\ref{fig:sweeps} with 95\% CIs).

\begin{table}[h]
\centering
\small
\begin{tabular}{c c c c c c}
\toprule
\textbf{collect\_noise} & \textbf{offline OPD} & \textbf{active query} & \textbf{full-dist adv} & \textbf{branch replay} & \textbf{Rang KD} \\
\midrule
0.00 & 0.673 & 0.693 & 0.699 & 0.699 & 0.699 \\
0.05 & 0.694 & 0.831 & 1.000 & 1.000 & 1.000 \\
0.15 & 0.601 & 0.853 & 1.000 & 1.000 & 1.000 \\
0.30 & 0.732 & 0.876 & 1.000 & 1.000 & 1.000 \\
\bottomrule
\end{tabular}
\caption{Success vs collection coverage, deploy fixed at 0.30. Active query benefits
from broader coverage (more states to query), while full-dist and branch replay hit
1.000 as soon as error states appear in data.}
\label{tab:coverage}
\end{table}

Plain Lightning OPD is flat across coverage ($\sim$0.6--0.73). Broadening
the dataset does not help, because the advantage gradient still ignores actions the
reference never sampled. Active query climbs with coverage (more states to query at),
but at the cost of env access. Both full-distribution methods sit at the SFT floor at
zero coverage, then jump to 1.000 the moment any noise surfaces the error states.
They only need the recovery states to appear, not to be acted out. Patch 2
tracks Rang cell for cell, and patch 3 matches them. That is the cleanest evidence
that what you distill is the operative variable.

\subsection{Ablation B --- deployment shift}
\code{results\_deploy.csv}, Figure~\ref{fig:deploy} (and panel (b) of
Figure~\ref{fig:sweeps} with 95\% CIs).

\begin{table}[h]
\centering
\small
\begin{tabular}{c c c c c c c}
\toprule
\textbf{deploy\_noise} & \textbf{offline OPD} & \textbf{full-dist adv} & \textbf{branch replay} & \textbf{Rang KD} & \textbf{online OPD} & \textbf{SFT} \\
\midrule
0.10 & 0.804 & 1.000 & 1.000 & 1.000 & 1.000 & 0.900 \\
0.20 & 0.751 & 1.000 & 1.000 & 1.000 & 1.000 & 0.795 \\
0.30 & 0.694 & 1.000 & 1.000 & 1.000 & 1.000 & 0.699 \\
0.45 & 0.608 & 1.000 & 1.000 & 1.000 & 1.000 & 0.554 \\
\bottomrule
\end{tabular}
\caption{Success vs deployment shift, collection fixed at 0.05. The three
full-distribution methods stay flat at 1.000 across all shift levels.}
\label{tab:deploy}
\end{table}

Lightning OPD drops steadily as the error branch fires more often, and it
tracks SFT closely. It inherits imitation's weakness under shift. The
full-distribution methods stay flat at 1.000. Shift only bites if the recovery
skill was never distilled.

\subsection{Ablation C --- generalization vs pure coverage}
\code{results\_representation.csv}, Figure~\ref{fig:repr}.

There are two ways an offline method can cover an error branch it never acted out.
Parameter generalization lets a shared-weight student interpolate on-path
behaviour into unseen states. Literal support means a per-state table only
knows states present in its data. To separate them I re-ran the offline methods
with a tabular per-state student, whose gradient at state $s$ touches only
$s$'s own logits. Off-support states get zero gradient and stay frozen at their SFT
value, by construction.

\begin{table}[h]
\centering
\small
\begin{tabular}{l c c c c}
\toprule
 & \textbf{collect 0.00} & \textbf{0.05} & \textbf{0.15} & \textbf{0.30} \\
\midrule
\textbf{linear} SFT          & 0.699 & 0.699 & 0.699 & 0.699 \\
\textbf{linear} Lightning OPD& 0.673 & 0.694 & 0.601 & 0.732 \\
\textbf{linear} Rang KD      & 0.699 & \textbf{1.000} & \textbf{1.000} & \textbf{1.000} \\
\textbf{tabular} SFT         & 0.697 & 0.697 & 0.697 & 0.697 \\
\textbf{tabular} Lightning OPD & 0.697 & 0.725 & 0.719 & 0.758 \\
\textbf{tabular} Rang KD     & 0.697 & 0.780 & 0.903 & 0.919 \\
online OPD (both)            & 1.000 & 1.000 & 1.000 & 1.000 \\
\bottomrule
\end{tabular}
\caption{Success at deploy 0.30 across the coverage sweep, linear vs tabular student.}
\label{tab:repr}
\end{table}

The tabular column makes two things clear.

First, at zero coverage every offline method collapses to exactly the SFT
floor (0.697). With no shared weights there is no gradient at the off-support
error states. Neither the advantage signal nor the full teacher distribution can
move the student there. The entire distillation gain over imitation, off-support,
is a generalization effect. Pull generalization out and Lightning OPD provably adds
nothing over SFT. Assumption 3.2 does not just degrade off-support, it leaves
the method with no signal at all. This disproves the ``capacity is the cure'' claim
(\S\ref{sec:assumptions}) directly. The tabular student is the highest-capacity
student in the study: one independent logit vector per state, zero approximation
error, the teacher exactly representable. The capacity reading predicts the easy
regime, yet at zero coverage it is the worst regime. Adding capacity this way makes
the result collapse to SFT, not approach online. Capacity without coverage of the
distilled signal buys exactly nothing.

Second, Rang's headline win was carried by generalization. In the linear
student, full-distribution KD jumps to 1.000 once recovery states appear, because
shared weights propagate that mass into the many error states the student reaches.
Strip generalization away and the same method, same data, tops out at $\sim$0.92.
Coverage does the work; generalization is the amplifier.

Online OPD stays at 1.000 in both representations, because it visits the
error states itself. That pins the blame: the tabular collapse is a coverage
failure, not a capacity or optimization one.

\begin{quote}\small
Note on framing. The off-support success rate does not crash to zero. The
on-path golden behaviour still solves the no-error episodes ($\sim$0.70). What
crashes to zero is the distillation gain over SFT off-support, the
quantity Assumption 3.2 is actually about.
\end{quote}

\subsection{When offline OPD works, and when it fails}

The experiments give a clean rule. Offline OPD works when the frozen
dataset already covers the states the student will visit, as in single-turn tasks, or
multi-turn tasks where deployment shift is small and the demos are diverse.
It fails when three things stack up: distribution shift moves the student
off the reference distribution, state coverage is incomplete there, and early
errors compound into states no demo ever reached. A multi-turn agentic task hits
all three at once. That is exactly where Lightning OPD lands below SFT.

% =====================================================================
\section{与 https://arxiv.org/pdf/2509.26497 中的离线策略相比有哪些区别？ \qtag{vi}}
\label{sec:rang}
% =====================================================================

The comparison paper is Rang et al., openPangu Embedded-1B
(\href{https://arxiv.org/pdf/2509.26497}{arXiv:2509.26497}). Both methods are
offline and distill a teacher. The difference is what they distill.

The two share the structure of the offline setting: no environment in the training
loop, a frozen teacher signal, and a student trained by gradient descent on a fixed
dataset. They differ in the target. Lightning OPD distills the sampled action's
advantage, $A_t = \log\pi_T(a_t) - \log\pi_\theta(a_t)$, on one action per state.
Rang distills the teacher's full next-token distribution, with a cross-entropy plus
KL loss, $\mathcal{L} = (1-\lambda)\,\mathrm{CE} + \lambda\,\mathrm{KL}(P\|Q)$,
$\lambda=0.7$. The full distribution carries the recovery action's probability mass
wherever the teacher puts it. The sampled-action advantage does not, because the
golden-path rollouts never sample \code{ask\_followup}.

This single difference is the whole story of Table~\ref{tab:metrics}. Rang reaches
1.000 from the same no-env budget. Lightning OPD does not.

\subsection{Is Rang's 100\% too good to be true? Stress-testing the headline}
\label{sec:rang-stress}

A zero-variance 1.000 should make a referee suspicious. So I tried to break it.

It is not 1.000 unconditionally. At \code{collect\_noise = 0}
(Table~\ref{tab:coverage}, and the linear Rang 0.00 cell of
Table~\ref{tab:repr}) Rang sits at 0.699, the SFT floor. A clean dataset never
surfaces the recovery states, so full-distribution KD has nothing to transfer. The
win depends on collection touching the error branch at least occasionally.

Strip generalization and the 1.000 disappears. The tabular column of
Table~\ref{tab:repr} is built for this. Same loss, same data, per-state student:
Rang tops out at 0.78 / 0.90 / 0.92 and never reaches 1.000. The headline was
generalization spending coverage efficiently, not the loss being complete.

I also measured where the recovery mass lives. On
golden-path states the teacher puts $\approx 10^{-8}$ on \code{ask\_followup},
effectively zero, not a small constant tail. The recovery mass (mean $\approx
0.6$, up to 1.0) sits entirely on the error-branch states. Those enter a frozen
dataset only when collection noise triggers them (0 such states at
\code{collect\_noise=0}; 32 / 56 / 99 records at 0.05 / 0.15 / 0.30). So Rang
wins because, once an error state appears, the full-distribution target there is
$\approx 1.0$ on recovery, while the sampled-action advantage at that same state
is whatever non-recovery action the rollout took.

So the 100\% is real but fragile and fully explained. It is the joint product of a
full-distribution target carrying recovery mass and parameter generalization
broadcasting it off-support. Remove the target and you get Lightning OPD (0.63).
Remove generalization and you get tabular Rang ($\le$0.92). Remove coverage and you
get the 0.699 floor. I report it as a ceiling conditional on its two enabling
mechanisms, not as an unqualified ``Rang wins.''

% =====================================================================
\section{针对观察到的某一个局限，提出一个代码/方法层面的改进，并做一个简单的消融实验 \qtag{vii}}
\label{sec:patch}
% =====================================================================

Lightning OPD evaluates only one sampled action per state, and the
offline dataset records that action. But the reference policy
($\pi_{\text{ref}}$) collected this data on a clean, optimal path. Once the agent
makes a mistake and drifts off-support, the recorded action is never the correct
recovery action. The reference policy does not know how to escape traps; its
probability of choosing the recovery action is near zero ($10^{-8}$). As a result,
the advantage formula cancels out the missing action entirely, and the final policy
gradient carries zero recovery signal. The student learns what not to do, but it
never learns how to save itself.

I propose three patches and ablate them. One breaks the method's premise; patches 2
and 3 are honest in-contract fixes. Patch 1, an uncertainty-triggered teacher query,
is SafeDAgger-style active imitation that queries the teacher only when
$\max_a \pi_\theta(a\mid s) < \tau$; it breaks the contract by spending the very
resource Lightning OPD exists to avoid, reaching greedy $0.880 \pm .13$ at $\tau=0.7$
with 1558 env-steps and 159 queries, and is analyzed in \S\ref{sec:patch1-contract}.
Patch 2, the full-distribution advantage, is the minimal in-contract fix: greedy 1.000,
sampled 0.867, 642 env-steps, zero queries, analyzed in \S\ref{sec:patch2}. Patch 3,
branch-aware replay, closes the remaining mass gap at zero env cost by reweighting
each record by $\alpha(s) = 1 + \gamma\cdot\mathrm{KL}(\pi_T \| \pi_{\mathrm{ref}})$,
upweighting error-branch states where teacher and reference maximally disagree; it
reaches greedy 1.000, sampled 0.902 ($\gamma=5$), up to 0.950 at $\gamma=100$, with
630 env-steps and zero queries, analyzed in \S\ref{sec:patch3}.

\subsection{Patch 1: strong but breaks the contract}
\label{sec:patch1-contract}

The professor's literal hint was to query the teacher only at uncertain states, a
SafeDAgger-style active imitation scheme. The intuition is sound: on the golden path
the student is confident, and only at error-branch states does it need help. I
implemented this as an uncertainty-triggered query: at each training step, if
$\max_a \pi_\theta(a\mid s) < \tau$ (confidence threshold), query the teacher for the
full distribution at that state and inject it into the dataset.

\begin{table}[h]
\centering
\small
\begin{tabular}{c c c c c}
\toprule
\textbf{conf thresh $\tau$} & \textbf{greedy} & \textbf{sampled} & \textbf{env-steps} & \textbf{queries} \\
\midrule
0.0  & $0.695 \pm .14$ & $0.711$ & 1594 & 0 \\
0.5  & $0.739 \pm .14$ & $0.744$ & 1586 & 36 \\
0.7  & $0.880 \pm .13$ & $0.826$ & 1558 & 159 \\
0.9  & $0.898 \pm .10$ & $0.853$ & 1560 & 191 \\
1.0  & $0.895 \pm .11$ & $0.847$ & 1581 & 931 \\
\bottomrule
\end{tabular}
\caption{Patch 1 ablation: confidence threshold $\tau$. Higher $\tau$ queries more
states, lifting greedy success but spending env-steps. Plotted with 95\% CIs in
panel (d) of Figure~\ref{fig:sweeps}.}
\label{tab:active-ablation}
\end{table}

The method works. At $\tau=0.7$, greedy climbs to $0.880$, a 0.25 lift over plain
Lightning OPD. But it has the same fatal flaw as the old DAgger patch:

\begin{quote}\small
Lightning OPD exists to remove the environment and the live teacher from
the training loop. That is the cost it sells the pots and pans to save. The moment
patch 1 calls \code{env.step} on student rollouts mid-training to reach the uncertain
states, it has spent exactly that resource back. And the 159 teacher queries are live
annotations on fresh states, another resource Lightning OPD freezes at collection.
\end{quote}

So patch 1 is not ``offline OPD, fixed.'' It is online learning in offline OPD's
clothes. Once you are paying for live env access and teacher queries, the honest
baseline is online OPD (1.000 greedy, 0.990 sampled), and patch 1's 0.880 is
strictly dominated within its own cost class. The professor's hint is a valid method
for online distillation with query budgeting. It is not a fix for Lightning
OPD's offline contract.

\begin{table}[h]
\centering
\small
\begin{tabular}{l c c c l}
\toprule
 & \textbf{env access?} & \textbf{live teacher?} & \textbf{greedy} & \textbf{honest verdict} \\
\midrule
Lightning OPD            & no  & no  & 0.63 & fails (the baseline) \\
active query (patch 1)   & \textbf{yes} & \textbf{yes} & 0.88 & \textbf{online learning in offline's clothes} \\
full-dist adv (patch 2)  & no  & no  & \textbf{1.00} & honest fix (full distribution) \\
branch replay (patch 3)  & no  & no  & \textbf{1.00} & honest fix (reweighted gradient) \\
Rang offline KD          & no  & no  & \textbf{1.00} & honest fix (full distribution) \\
online OPD               & yes & yes & 1.00 & the right tool once env is on the table \\
online RL                & yes & no  & 1.00 & the right tool once env is on the table \\
\bottomrule
\end{tabular}
\caption{The honest frontier. No row is both no-env-access and beats the
full-distribution offline methods. Patch 1 sits strictly inside the frontier.}
\label{tab:frontier}
\end{table}

\begin{table}[h]
\centering
\small
\begin{tabular}{l c r l c}
\toprule
\textbf{method} & \textbf{env access} & \textbf{env-steps} & \textbf{spent\ldots} & \textbf{greedy} \\
\midrule
teacher / SFT             & none & \textbf{0}     & ---                            & 1.000 / 0.699 \\
\textbf{Lightning OPD}    & none & \textbf{633}   & once, at collection            & 0.632 \\
\textbf{branch replay (patch 3)} & none & \textbf{630} & once, at collection       & \textbf{1.000} \\
\textbf{full-dist adv (patch 2)} & none & \textbf{635} & once, at collection       & \textbf{1.000} \\
Rang offline KD           & none & 642            & once, at collection            & \textbf{1.000} \\
active query (patch 1)    & \textbf{yes} & \textbf{1\,558} & once + \textbf{mid-training rollouts} & 0.880 \\
online OPD                & yes  & \textbf{30\,303}  & \textbf{every gradient step} & 1.000 \\
online RL                 & yes  & \textbf{154\,526} & \textbf{every gradient step} & 1.000 \\
\bottomrule
\end{tabular}
\caption{Measured environment budget. The timing column is the honest axis.
The three no-env patches share Lightning OPD's $\approx$630 budget; patch 1 pays
2.5$\times$ more for weaker results.}
\label{tab:budget}
\end{table}

Two things these tables make clear. First, the four no-env methods share one budget:
Lightning OPD, patches 2 and 3, and Rang all pay $\approx$ 630--642 env-steps, a
one-time collection that training reuses with zero further \code{env.step}. On that
identical budget the score ranges from 0.632 to 1.000, so the comparison is
apples-to-apples, and what separates the methods is the loss, not the budget. Second,
patch 1 and the online methods are a different cost class. Patch 1's 1\,558 steps are
recurring mid-training rollouts plus teacher queries, and online OPD/RL spend
env-steps on every gradient step, 48$\times$ / 244$\times$ the offline budget. Once a
method is in this column, the honest baseline is online OPD/RL (both 1.000), and patch
1's 0.880 is strictly dominated within its own cost class.

Figure~\ref{fig:pareto} draws this frontier on its true axis. Plotting greedy success
against total interaction budget makes the apples-to-apples comparison visible at a
glance: the four no-env methods cluster at $\approx$630 steps, the in-contract patches
and Rang ride the same 1.000 ceiling as online OPD/RL at $\sim$50$\times$ less
interaction, and patch~1 sits strictly inside the frontier.

\subsection{Patch 2: the minimal in-contract fix (full-distribution advantage)}
\label{sec:patch2}

Rang's 1.000 leaves an objection: maybe it wins because it is a different method
(CE+KL), not because of the full-distribution target specifically. Patch 2 is the
control that isolates the variable. It keeps Lightning OPD's gradient
exactly, with the same frozen dataset, no env access, no live teacher, the same
advantage-PG form, and the same clip, and changes one thing. Instead of the sampled
action's advantage $A(a)\cdot\nabla\log\pi(a)$, it takes the expectation under the
teacher's full distribution:
\[
\sum_a \pi_T(a\mid s)\,\cdot\,\mathrm{clip}\bigl(\log\pi_T(a\mid s) - \log\pi_\theta(a\mid s)\bigr)\,\cdot\,\nabla\log\pi_\theta(a\mid s).
\]
That is the single edit. The result is $1.000 \pm .000$, tying Rang and tracking it
cell for cell across the coverage sweep. It runs at the same $\sim$635 env-step
budget as plain Lightning OPD, with zero additional interaction.

\paragraph{Why this is not a variance argument.}
It is tempting to read patch 2 as ``the same gradient, lower variance.'' That is
wrong, and the distinction is the crux of the diagnosis. Write the per-state
advantage-PG contribution. Lightning OPD evaluates it at the single rollout action
$a \sim b(\cdot\mid s)$, where $b$ is the behaviour policy that generated the
dataset (here $\pi_{\mathrm{ref}}$):\[
g_{\mathrm{LOPD}}(s) = A(a,s)\,\nabla\log\pi_\theta(a\mid s),
\qquad a \sim b(\cdot\mid s).
\]
Patch 2 replaces the sampled action by the expectation under the teacher:
\[
g_{\mathrm{FULL}}(s) = \sum_a \pi_T(a\mid s)\,A(a,s)\,\nabla\log\pi_\theta(a\mid s)
= \mathbb{E}_{a \sim \pi_T(\cdot\mid s)}\!\bigl[A(a,s)\,\nabla\log\pi_\theta(a\mid s)\bigr].
\]
The catch is the sampling distribution. A Monte-Carlo estimate of $g_{\mathrm{FULL}}$
needs $a \sim \pi_T$, but Lightning OPD's one sample is drawn from
$b = \pi_{\mathrm{ref}}$. So $g_{\mathrm{LOPD}}$ is an unbiased one-sample estimator
of a different quantity,
\[
\mathbb{E}_{a \sim \pi_{\mathrm{ref}}}\!\bigl[A(a,s)\,\nabla\log\pi_\theta(a\mid s)\bigr]
\;\neq\;
\mathbb{E}_{a \sim \pi_T}\!\bigl[A(a,s)\,\nabla\log\pi_\theta(a\mid s)\bigr]
= g_{\mathrm{FULL}}(s),
\]
and the two coincide only when $\pi_{\mathrm{ref}} \approx \pi_T$ on
$\mathrm{supp}(\pi_T(\cdot\mid s))$. That is exactly the support-coverage
assumption, re-derived as an importance-sampling condition. Off-support it fails
maximally. At an error-branch state the recovery action $r=\code{ask\_followup}$ has
$\pi_T(r\mid s)\approx 1.0$ but $\pi_{\mathrm{ref}}(r\mid s)\approx 0$ (measured
$\approx 10^{-8}$). The term that dominates $g_{\mathrm{FULL}}$,
\[
\pi_T(r\mid s)\,A(r,s)\,\nabla\log\pi_\theta(r\mid s)
\quad\text{with}\quad \pi_T(r\mid s)\approx 1.0,\ A(r,s)\gg 0,
\]
is sampled by $g_{\mathrm{LOPD}}$ with probability $\pi_{\mathrm{ref}}(r\mid s)
\approx 0$. No amount of averaging over more rollouts recovers it, because
every rollout is drawn from the same $\pi_{\mathrm{ref}}$ that puts zero mass on
$r$. The missing term is not high-variance. It is systematically absent from the
estimator's support. Increasing the dataset is the wrong axis, since it shrinks
variance on terms $\pi_{\mathrm{ref}}$ already samples and does nothing for the term
it never can. This is the formal content of ``flat across coverage'' in
Table~\ref{tab:coverage}.

Two clarifications keep this honest. First, $g_{\mathrm{FULL}}$ is not
the gradient of $\mathrm{KL}(\pi_T\|\pi_\theta)$, since that gradient is
$\sum_a \pi_T(a)\nabla\log\pi_\theta(a)$, with no advantage weight.
$g_{\mathrm{FULL}}$ is the advantage-weighted policy gradient with the teacher as
sampling measure. It keeps Lightning OPD's stop-grad advantage word for word and changes
only the measure, from $\pi_{\mathrm{ref}}$ (one sample) to $\pi_T$ (closed form).
What the two share is the fixed point: at $\pi_\theta = \pi_T$ every $A(a,s)=0$, so
both vanish. Second, patch 2 is therefore not a variance-reduced estimator of
$g_{\mathrm{LOPD}}$'s target. It computes a different target exactly, using
the one piece of information Lightning OPD already precomputed but only ever read at
the sampled index, the full row $\pi_T(\cdot\mid s)$. It costs no env access
because the recovery mass lives in $\pi_T$ (precomputed offline), not in
$\pi_{\mathrm{ref}}$ (which did the sampling).

This pins the diagnosis with no room left to argue. It is not the no-env-access
contract that breaks Lightning OPD, since patch 2 keeps it and hits ceiling. It is not
that Rang's CE+KL loss is magic, since a one-line change to Lightning OPD's own gradient
reproduces the win. It is the sampled-action restriction: formally, estimating
$g_{\mathrm{FULL}}$ by a single draw from $\pi_{\mathrm{ref}}$ instead of
evaluating the $\pi_T$-expectation in closed form.

Patch 2 is the honest fix the paper should have shipped. It stays inside the exact
contract Lightning OPD sells, and still closes the gap, because the gap was never
about the contract. It was about distilling one action instead of the whole
distribution.

\subsection{Patch 3: closing the mass gap at zero env cost (branch-aware replay)}
\label{sec:patch3}

Patch 2 hits greedy 1.000 but drops to 0.867 under sampling, a $-0.13$ gap. The
recovery action is the mode at error states, but $\sim$13\% of the mass still
sits on wrong actions. Online OPD closes this gap (sampled 0.990) by visiting
the error states and reinforcing recovery to dominance. Patch 3 asks: can we close
the mass gap without env access?

The idea is branch-aware replay. The frozen dataset already contains
error-branch records, which appeared when collection noise drifted the reference
off the golden path. But they are rare (off-support ratio $\approx 0.08$), so their
gradient contribution is diluted. Patch 3 reweights each record by
\[
\alpha(s) = 1 + \gamma\cdot\mathrm{KL}\bigl(\pi_T(\cdot\mid s) \| \pi_{\mathrm{ref}}(\cdot\mid s)\bigr).
\]
At error-branch states, the teacher strongly prefers \code{ask\_followup}
($\pi_T \approx 1.0$) while the reference never learned it
($\pi_{\mathrm{ref}} \approx 10^{-8}$). The KL is large, so $\alpha(s)$ upweights
these rare but critical records. On the golden path, $\pi_T \approx \pi_{\mathrm{ref}}$,
so $\alpha(s) \approx 1$ and records pass through unchanged.

\begin{table}[h]
\centering
\small
\begin{tabular}{c c c c c}
\toprule
\textbf{$\gamma$} & \textbf{greedy} & \textbf{sampled} & \textbf{env-steps} & \textbf{queries} \\
\midrule
0.0   & 1.000 & 0.870 & 626 & 0 \\
1.0   & 1.000 & 0.882 & 627 & 0 \\
5.0   & 1.000 & 0.902 & 630 & 0 \\
20.0  & 1.000 & 0.932 & 627 & 0 \\
100.0 & 1.000 & 0.950 & 627 & 0 \\
\bottomrule
\end{tabular}
\caption{Patch 3 ablation: KL weight $\gamma$. Higher $\gamma$ lifts sampled success
with zero env cost, closing $\sim$80\% of the greedy--sampled gap at $\gamma=100$.
Plotted with 95\% CIs in panel (c) of Figure~\ref{fig:sweeps}.}
\label{tab:branch-ablation}
\end{table}

The ablation tells the story. At $\gamma=0$ (uniform weighting) we recover patch 2's
0.870 sampled baseline. As $\gamma$ grows, sampled success climbs monotonically:
0.882 at $\gamma=1$, 0.902 at $\gamma=5$, 0.932 at $\gamma=20$, reaching 0.950 at
$\gamma=100$. That closes $\sim$80\% of the 0.13 gap between greedy 1.000 and online
OPD's sampled 0.990, all at the same $\approx$630 env-step budget, with zero additional
interaction and zero teacher queries.

\paragraph{Why this works without env access.}
The error-branch records already exist in the frozen dataset; collection noise
surfaced them. The problem is not absence but dilution: they contribute
$\approx$8\% of the gradient budget while the golden-path records dominate at 92\%.
Patch 3 does not create new data. It reweights existing data so the rare error-branch
records pull their weight. The KL between $\pi_T$ and $\pi_{\mathrm{ref}}$ is computed
from the frozen teacher distribution and the frozen reference policy, both
precomputed at collection, with no live annotation needed.

\paragraph{What remains.}
Even at $\gamma=100$, sampled 0.950 trails online OPD's 0.990 by 0.04. That residual is
the true cost of never visiting the error states. Online methods reinforce
recovery mass to $\approx$1.0 by receiving reward signals at the states. Offline
methods only distill what the teacher would do, never seeing the outcome.
Patch 3 squeezes the most out of frozen data, but the last 0.04 is exactly what
env access buys.

\subsection{What the 1.000 hides: greedy vs sampled evaluation}
\label{sec:greedy}

The headline metric takes greedy (argmax) actions. That is standard, but it
makes offline distillation look better than it is. A careful referee should ask what distribution the
student actually learned, not just its mode. So I re-evaluated every method by
sampling from the policy (\code{sampled\_success\_rate}).

\begin{table}[h]
\centering
\small
\begin{tabular}{l c c c c}
\toprule
\textbf{method} & \textbf{greedy} & \textbf{sampled} & \textbf{$\Delta$ (sampled $-$ greedy)} & \textbf{env access} \\
\midrule
teacher       & 1.000 & 1.000 & $0.00$  & --- \\
online RL     & 1.000 & 0.997 & $-0.00$ & yes \\
online OPD    & 1.000 & 0.990 & $-0.01$ & yes \\
\textbf{branch replay (patch 3)} & \textbf{1.000} & \textbf{0.902} & $\mathbf{-0.10}$ & no \\
\textbf{Rang offline KD} & \textbf{1.000} & \textbf{0.869} & $\mathbf{-0.13}$ & no \\
\textbf{full-dist adv (patch 2)} & \textbf{1.000} & \textbf{0.867} & $\mathbf{-0.13}$ & no \\
active query (patch 1) & 0.880 & 0.826 & $-0.05$ & yes \\
SFT           & 0.699 & 0.749 & $+0.05$ & no \\
Lightning OPD & 0.632 & 0.665 & $+0.03$ & no \\
\bottomrule
\end{tabular}
\caption{Greedy vs sampled success. Patch 3 closes $\sim$80\% of the greedy--sampled
gap that patch 2 leaves, reaching 0.902 sampled at zero env cost.}
\label{tab:greedy}
\end{table}

Two honest admissions, both of which sharpen the thesis.

My own 1.000s are greedy-conditional. The three no-env offline methods
that hit greedy 1.000 all drop under sampling: patch 2 to 0.867 ($-0.13$), Rang to
0.869 ($-0.13$), and, crucially, patch 3 to 0.902 ($-0.10$). So even the best
offline fix does not learn a confidently recovering policy. It learns one where
\code{ask\_followup} is merely the most likely action at the error states, with
$\sim$10--13\% of the mass still on wrong actions. Patch 3's branch-aware reweighting
closes $\sim$80\% of this gap at $\gamma=100$ (sampled 0.950), but the residual 0.04
is exactly what env access buys.

The greedy--sampled gap is itself a coverage thermometer. The two
online methods barely move ($-0.00/-0.01$): having visited and been rewarded at
the error states, they drive recovery mass to $\approx 1.0$, so sampling and
argmax agree. The offline methods lose 0.10--0.13 precisely because off-support they
only ever saw the recovery action once it appeared in data, never reinforced it to
dominance. Patch 3 narrows the drop to 0.10 by reweighting error-branch records,
but without env access it cannot reach online's 0.01. The size of the drop separates
having visited the states (online, $\sim$0) from having only distilled the states
(patch 2, $\sim$0.13) from distilling with branch-aware replay (patch 3, $\sim$0.10).

So even my best fix carries the fingerprint of the disease. The honest framing is not
``patches 2 and 3 solve it'' but ``patch 2 recovers the mode off-support at zero env
cost, patch 3 closes 80\% of the remaining mass gap still at zero env cost, and the
last 0.04 is exactly what costs env access (online OPD/RL, sampled $\approx 0.99$).''

% =====================================================================
\section{Takeaway}
\label{sec:takeaway}
% =====================================================================

I do not attack Lightning OPD's admitted open problem (multi-turn coverage). I
attack its central, load-bearing claim: that under teacher consistency the
online--offline gap is governed by drift (Thm 3.5), and any residual is
capacity-limited. I satisfy that claim's premises exactly. Consistency holds by
construction ($\sigma_\Delta^2 = 0$). Capacity is sufficient (the same linear
student reaches 1.000 once fed the full distribution). And the method still lands at
or below SFT (0.632 vs 0.699), with an order of magnitude more variance. The
surviving gap is therefore neither a consistency bias nor a capacity limit. It is a
coverage-of-the-distilled-signal failure. Three messages, in order of force.

First, the failure is not where the paper's theory locates it. With
consistency satisfied and capacity sufficient, the method still fails, because
off-support it has no signal, and only literal support or live env access can
supply one. The ``capacity is the cure'' claim is proven wrong by the
tabular column, where maximal capacity at zero coverage is the worst case.

Second, given generalization, what you distill decides whether
coverage is enough. Patch 2 isolates the variable: hold the dataset, the no-env
contract, and the advantage-PG form fixed, swap only the sampled action for the
full distribution, and the score goes $0.63 \to 1.000$, one line, in-contract.
Under sampling that 1.000 is $\approx 0.87$, so patch 2 recovers the mode, not the
full mass. Patch 3 closes 80\% of that gap at $\gamma=100$, still in-contract.

Third, the results that look like they cut against the thesis don't. Active query's
0.88 is online learning in disguise: it spends the env access and teacher queries
Lightning OPD exists to avoid, and is dominated by the online methods it has become.
Rang's 1.000 is not unconditional: it falls to the SFT floor at zero coverage and
to $\le$0.92 once generalization is removed. Every surprise reduces to the same
mechanism as the main failure.

% =====================================================================
\section*{Reproduce}
% =====================================================================

\begin{verbatim}
cd opd_toy
python run_experiments.py     # writes ../results/*.csv + case_studies.txt
python plots.py               # writes ../results/fig_*.png
python plots_paper.py         # writes ../results/fig_{pareto,sweeps}.{pdf,png}
\end{verbatim}

Artifacts: \code{results\_main.csv}, \code{results\_coverage.csv},
\code{results\_deploy.csv}, \code{results\_representation.csv},
\code{results\_active.csv}, \code{results\_patch3.csv}, \code{case\_studies.txt}, and
\code{fig\_\{main,coverage,deploy,representation,pareto,sweeps\}.png}.

% =====================================================================
% Figures
% =====================================================================

\begin{figure}[h]
\centering
\includegraphics[width=0.85\textwidth]{fig_main.png}
\caption{Headline success rate at the fixed operating point (5 seeds). Lightning
OPD lands below SFT with the largest variance in the table.}
\label{fig:main}
\end{figure}

\begin{figure}[h]
\centering
\includegraphics[width=0.85\textwidth]{fig_coverage.png}
\caption{Ablation A. Plain Lightning OPD is flat across collection coverage; the
full-distribution methods jump to 1.000 once error states appear.}
\label{fig:coverage}
\end{figure}

\begin{figure}[h]
\centering
\includegraphics[width=0.85\textwidth]{fig_deploy.png}
\caption{Ablation B. Lightning OPD drops steadily under deployment shift,
tracking SFT; the full-distribution methods stay flat at 1.000.}
\label{fig:deploy}
\end{figure}

\begin{figure}[h]
\centering
\includegraphics[width=0.85\textwidth]{fig_representation.png}
\caption{Ablation C. Linear vs tabular student. At zero coverage every offline
method collapses to the SFT floor.}
\label{fig:repr}
\end{figure}

\begin{figure}[h]
\centering
\includegraphics[width=0.92\textwidth]{fig_pareto.png}
\caption{The honest frontier, drawn on its true axis. Greedy success rate vs.\
total interaction budget (env-steps $+$ teacher queries, symlog $x$), 5 seeds with
95\% Student-$t$ confidence intervals. The two in-contract patches (full-distribution
advantage, branch-aware replay) and Rang reach the same 1.000 ceiling as online
OPD/RL at $\sim$50$\times$ less interaction, sitting on the dashed Pareto frontier.
Active query (patch~1, hollow marker) breaks the no-env/no-teacher contract and is
dominated within its own cost class. Plain Lightning OPD carries the widest interval,
reflecting its severe seed variance.}
\label{fig:pareto}
\end{figure}

\begin{figure}[h]
\centering
\includegraphics[width=0.95\textwidth]{fig_sweeps.png}
\caption{Consolidated ablations, 5 seeds with 95\% Student-$t$ confidence bands.
(a) Collection coverage: plain Lightning OPD stays flat and noisy while the
full-distribution methods jump to 1.000 once error states appear.
(b) Deployment shift: Lightning OPD tracks the SFT floor downward; full-distribution
adv.\ stays flat near the online ceiling.
(c) Branch-replay KL weight $\gamma$ (symlog $x$): greedy is saturated, so sampled
success is plotted; it climbs monotonically at zero env cost.
(d) Active-query confidence threshold $\tau$: success rate (left axis) rises with the
teacher-query count (right axis), exposing the accuracy-for-interaction trade.}
\label{fig:sweeps}
\end{figure}

\end{document}
